%
% Teil 2: Auswertung des Bromwich-Integrals mit dem Residuensatz
% und dem Lemma von Jordan (Schliessen der Kontur ueber den im Ursprung
% zentrierten Kreisbogen, D-Kontur)
%
\section{Auswertung mit dem Residuensatz: die geschlossene D-Kontur}
\label{laplace:subsection:residuensatz}

Mit Gleichung~\eqref{laplace:eq:bromwich_integral} ist die Inversion zwar
formal gelöst, praktisch stehen wir aber vor einem unendlich langen
Integrationsweg mitten in der komplexen Ebene.
Der Schlüssel zur konkreten Auswertung ist der Residuensatz, der im ersten
Teil dieses Buches ausführlich hergeleitet wurde und dessen Anwendung auf
Konturintegrale beispielsweise in \cite{laplace:ablowitz2003} vertieft wird.
Ist $G(s)$ innerhalb einer geschlossenen, positiv orientierten Kurve $C$
holomorph bis auf endlich viele isolierte Singularitäten $s_1,\dots,s_n$,
so gilt
\begin{equation}
    \oint_C G(s) \, ds
    =
    2\pi i \sum_{k=1}^{n} \operatorname{Res}\bigl(G(s), s_k\bigr).
    \label{laplace:eq:residuensatz}
\end{equation}
Entscheidend ist dabei ein oft übersehenes Detail:
Nur die von der Kurve \emph{eingeschlossenen} Singularitäten tragen zum
Integral bei.
Eine Polstelle, die auch nur knapp ausserhalb der Kontur liegt, ist für das
Integral schlicht unsichtbar.
Abbildung~\ref{laplace:fig:pol_auswahl} veranschaulicht diese Auswahlregel.

\begin{figure}
    \centering
    \scalebox{0.85}{
\begin{tikzpicture}

% ==============================
% 1. FARBEN DEFINIEREN
% ==============================
\definecolor{plotblue}{RGB}{0, 0, 200}
\definecolor{plotred}{RGB}{220, 30, 30}
\definecolor{plotgray}{RGB}{140, 140, 140}
\definecolor{plotbg}{RGB}{244, 244, 252}

% ==============================
% 2. ACHSEN
% ==============================
\draw[thick, -{Stealth[length=3.5mm, width=2.5mm]}] (-4.2, 0) -- (4.5, 0) node[right] {\textbf{Re}$(s)$};
\draw[thick, -{Stealth[length=3.5mm, width=2.5mm]}] (0, -4.2) -- (0, 4.2) node[above] {\textbf{Im}$(s)$};

% ==============================
% 3. KONTUR C
% ==============================
% Geschlossene, organische Blob-Form, positiv orientiert
\draw[thick, plotblue, decoration={markings, mark=at position 0.25 with {\arrow{Stealth[length=3mm, width=2mm]}},
                          mark=at position 0.60 with {\arrow{Stealth[length=3mm, width=2mm]}},
                          mark=at position 0.90 with {\arrow{Stealth[length=3mm, width=2mm]}}}, postaction={decorate}] 
    (-0.5, 0) .. controls (1, 3) and (3, 1) .. 
    (2, -1) .. controls (1.5, -3) and (-3, -3) .. 
    (-2.5, -1) .. controls (-2, 1) and (-1, -1.5) .. (-0.5, 0);

% Label C aussen
\node[plotblue, font=\large] at (1.8, 1.8) {$C$};

% ==============================
% 4. POLE (Kreuze)
% ==============================
\newcommand{\pole}[3][plotred]{
    \draw[thick, #1, line cap=round] (#2-0.12,#3-0.12) -- (#2+0.12,#3+0.12);
    \draw[thick, #1, line cap=round] (#2-0.12,#3+0.12) -- (#2+0.12,#3-0.12);
}

% Innere Pole (rot)
\pole{0.8}{0.8}
\pole{0.0}{-1.8}
\pole{-1.0}{-1.2}

% Aeussere Pole (grau)
\pole[plotgray]{3.0}{2.0}
\pole[plotgray]{-2.8}{-2.0}

% ==============================
% 5. LEGENDE
% ==============================
\node[draw, thick, rounded corners=3pt, fill=white, inner xsep=5pt, inner ysep=4pt, anchor=north west] at (-4.0, 4.0) {
    \renewcommand{\arraystretch}{1.2}
    \begin{tabular}{@{}c l@{}}
        
        \tikz[baseline=-0.6ex]{
            \draw[thick, plotred, line cap=round] (-0.1,-0.1) -- (0.1,0.1) (-0.1,0.1) -- (0.1,-0.1);
        } & Relevant \\
        
        \tikz[baseline=-0.6ex]{
            \draw[thick, plotgray, line cap=round] (-0.1,-0.1) -- (0.1,0.1) (-0.1,0.1) -- (0.1,-0.1);
        } & Irrelevant \\
        
        \tikz[baseline=-0.6ex]{
            \draw[thick, plotblue, ->-=0.6] (-0.3, 0) -- (0.3, 0);
        } & Kontur $C$ \\
        
    \end{tabular}
};

\end{tikzpicture}
}

    \caption{Auswahlregel des Residuensatzes:
    Nur die von der geschlossenen Kontur $C$ eingeschlossenen Polstellen
    (rot) liefern einen Beitrag $2\pi i \operatorname{Res}$ zum
    Konturintegral.
    Die Polstelle ausserhalb der Kontur (grau) ist für das Integral
    unsichtbar.}
    \label{laplace:fig:pol_auswahl}
\end{figure}

Der Residuensatz verlangt allerdings eine \emph{geschlossene} Kurve, während
der Bromwich-Pfad eine offene, unendlich lange Gerade ist.
Wir müssen die Gerade also künstlich schliessen, und zwar so, dass der
hinzugefügte Wegabschnitt im Grenzfall nichts zum Integral beiträgt.
Die Richtung des Schliessens gibt der Faktor $e^{st}$ vor.
Wegen
\begin{equation}
    \bigl| e^{st} \bigr|
    =
    e^{t \, \text{Re}(s)}
    \label{laplace:eq:betrag_est}
\end{equation}
fällt der Integrand für $t > 0$ exponentiell ab, sobald wir uns in die linke
Halbebene $\text{Re}(s) < 0$ begeben, und er explodiert in der rechten.
Wir schliessen die Kontur daher nach links, und zwar über einen grossen,
im Ursprung zentrierten Kreisbogen $\Gamma_R$, wie ihn auch
\cite{laplace:carslaw1959} verwendet.
Da sämtliche Singularitäten von $F(s)$ links der Geraden
$\text{Re}(s) = \gamma$ liegen, fängt die geschlossene Kurve sie für
hinreichend grosses $R$ alle ein.
Zusammen mit dem vertikalen Geradenstück, das als Sehne des Kreises
zwischen den beiden Schnittpunkten mit dem Bogen verläuft, entsteht die
in Abbildung~\ref{laplace:fig:d_kontur} gezeigte geschlossene Kurve, die
wegen ihrer Form als \emph{D-Kontur} bezeichnet wird.
Dass der ursprungszentrierte Bogen dabei im Streifen
$0 < \text{Re}(s) < \gamma$ ein kleines Stück in die rechte Halbebene
hineinragt, wird sich gleich als harmlos erweisen.

\begin{figure}
    \centering
    \scalebox{0.85}{
\begin{tikzpicture}

% ==============================
% 1. FARBEN DEFINIEREN
% ==============================
\definecolor{plotblue}{RGB}{0, 0, 200}
\definecolor{plotred}{RGB}{220, 30, 30}
\definecolor{plotgray}{RGB}{140, 140, 140}
\definecolor{plotorange}{RGB}{240, 140, 0}
\definecolor{plotbg}{RGB}{244, 244, 252}

% ==============================
% 2. HINTERGRUND FUELLEN
% ==============================
%\fill[plotbg] (1.5, -3.923) -- (1.5, 3.923) arc (69.07:290.93:4.2) -- cycle;

% ==============================
% 3. ACHSEN
% ==============================
\draw[thick, -{Stealth[length=3.5mm, width=2.5mm]}] (-5.5, 0) -- (3.5, 0) node[right] {\textbf{Re}$(s)$};
\draw[thick, -{Stealth[length=3.5mm, width=2.5mm]}] (0, -5.0) -- (0, 5.5) node[above] {\textbf{Im}$(s)$};

% ==============================
% 4. D-KONTUR
% ==============================

% Bromwich-Pfad (Sehne)
\draw[thick, plotblue, decoration={markings, mark=at position 0.55 with {\arrow{Stealth[length=3mm, width=2mm]}}}, postaction={decorate}] 
    (1.5, -3.923) -- (1.5, 3.923);

% Gamma Tick
\draw[thick] (1.5, 0.1) -- (1.5, -0.1) node[below right, inner sep=2pt] {$\boldsymbol{\gamma}$};

% Ticks an den Enden des Bromwich-Pfades
% (Schnittpunkte Sehne/Kreis pythagoreisch exakt: gamma +- i*sqrt(R^2-gamma^2);
%  fuer R -> infty naehern sie sich gamma +- iR an)
\draw[thick] (1.35, 3.923) -- (1.65, 3.923) node[right, text=black] {$\gamma + i\sqrt{R^2 - \gamma^2}$};
\draw[thick] (1.35, -3.923) -- (1.65, -3.923) node[right, text=black] {$\gamma - i\sqrt{R^2 - \gamma^2}$};

% Kreisbogen Gamma_R
\draw[thick, plotorange, decoration={markings, 
        mark=at position 0.3 with {\arrow{Stealth[length=3mm, width=2mm]}},
        mark=at position 0.7 with {\arrow{Stealth[length=3mm, width=2mm]}}}, 
        postaction={decorate}] 
    (1.5, 3.923) arc (69.07:290.93:4.2);

\node[plotorange, font=\large] at (-4.6, 1.8) {$\Gamma_R$};

% Markierung fuer den Bereich in der rechten Halbebene
% Ragt in rechte Halbebene (entfernt)

% Radiuspfeil
\draw[thick, plotgray, dashed, -{Stealth[length=2.5mm, width=2mm]}] (0,0) -- (-2.97, 2.97) node[midway, below left] {$R\to\infty$};

% ==============================
% 5. POLE (Kreuze)
% ==============================
\newcommand{\pole}[3][plotred]{
    \draw[thick, #1, line cap=round] (#2-0.12,#3-0.12) -- (#2+0.12,#3+0.12);
    \draw[thick, #1, line cap=round] (#2-0.12,#3+0.12) -- (#2+0.12,#3-0.12);
}

% Innere Pole (rot)
\pole{0.0}{2.0}
\pole{0.0}{-2.0}
\pole{-1.2}{0.9}
\pole{-2.0}{-0.5}
\pole{-0.5}{0.0}

% ==============================
% 6. LEGENDE
% ==============================
\node[draw, thick, rounded corners=3pt, fill=white, inner xsep=5pt, inner ysep=4pt, anchor=north west] at (-5.3, 5.3) {
    \renewcommand{\arraystretch}{1.2}
    \begin{tabular}{@{}c l@{}}
        
        \tikz[baseline=-0.6ex]{
            \draw[thick, plotred, line cap=round] (-0.1,-0.1) -- (0.1,0.1) (-0.1,0.1) -- (0.1,-0.1);
        } & Polstelle \\
        
        \tikz[baseline=-0.6ex]{
            \draw[thick, plotblue, decoration={markings, mark=at position 0.55 with {\arrow{Stealth[length=2mm, width=1.5mm]}}}, postaction={decorate}] (-0.3, 0) -- (0.3, 0);
        } & Bromwich-Pfad \\
        
        \tikz[baseline=-0.6ex]{
            \draw[thick, plotorange, decoration={markings, mark=at position 0.55 with {\arrow{Stealth[length=2mm, width=1.5mm]}}}, postaction={decorate}] (-0.3, 0) -- (0.3, 0);
        } & Kreisbogen $\Gamma_R$ \\
        
    \end{tabular}
};

\end{tikzpicture}
}

    \caption{Die geschlossene D-Kontur:
    Der vertikale Bromwich-Pfad bildet die Sehne eines im Ursprung
    zentrierten Kreises, dessen grosser Bogen $\Gamma_R$ die Kurve nach
    links schliesst und dabei im Streifen $0 < \text{Re}(s) < \gamma$
    leicht in die rechte Halbebene hineinragt.
    Für $R \to \infty$ umschliesst die Kurve alle Singularitäten von
    $F(s)$, während der Beitrag des Bogens nach dem Lemma von Jordan
    verschwindet.}
    \label{laplace:fig:d_kontur}
\end{figure}

Dass der Bogenbeitrag tatsächlich verschwindet, ist keineswegs
selbstverständlich, denn die Bogenlänge wächst proportional zu $R$ ins
Unendliche.
Tief in der linken Halbebene ist das unkritisch:
Dort erdrückt der exponentielle Abfall von $e^{st}$ aus
Gleichung~\eqref{laplace:eq:betrag_est} das nur lineare Anwachsen der
Weglänge mühelos.
Kritisch sind allein die Bogenstücke nahe der imaginären Achse und im
schmalen Streifen bis $\text{Re}(s) = \gamma$, wo
$|e^{st}| \leq e^{\gamma t}$ lediglich beschränkt ist; dort muss der
gleichmässige Abfall der Bildfunktion einspringen.
Genau das leistet das \emph{Lemma von Jordan} in der für die
Laplace-Inversion passenden Form, wie sie etwa in
\cite{laplace:doetsch1961} als Satz 34.2 bewiesen wird:
Strebt $F(s)$ auf den Bögen für $|s| \to \infty$ gleichmässig gegen
null, so gilt für $t > 0$
\begin{equation}
    \lim_{R \to \infty} \int_{\Gamma_R} e^{st} F(s) \, ds
    =
    0.
    \label{laplace:eq:jordan}
\end{equation}
Das Resultat deckt auch die rechts der imaginären Achse liegenden
Bogenstücke ab, denn deren geometrische Länge bleibt für $R \to \infty$
beschränkt:
Jedes der beiden Bogenstücke durchquert den Streifen
$0 < \text{Re}(s) < \gamma$ nahezu senkrecht, und seine Länge
$R \arcsin(\gamma/R)$ strebt für $R \to \infty$ gegen den festen Wert
$\gamma$.
Auf einem Wegstück beschränkter Länge, auf dem $|e^{st}| \leq e^{\gamma t}$
gilt, löscht der gleichmässige Abfall von $F(s)$ den Beitrag ebenfalls
aus.
Für die in der Praxis auftretenden Bildfunktionen, etwa gebrochen rationale
Funktionen mit Nennergrad grösser als Zählergrad, ist diese Voraussetzung
stets erfüllt.

Damit fällt das Puzzle zusammen.
Lassen wir in der D-Kontur den Radius $R$ gegen unendlich streben, so
verschwindet der Bogenbeitrag nach
Gleichung~\eqref{laplace:eq:jordan}, das vertikale Geradenstück wird zum
Bromwich-Integral, und der Residuensatz
\eqref{laplace:eq:residuensatz} liefert die bemerkenswert kompakte
Inversionsformel
\begin{equation}
    f(t)
    =
    \frac{1}{2\pi i} \int_{\gamma - i\infty}^{\gamma + i\infty}
    F(s) e^{st} \, ds
    =
    \sum_{k} \operatorname{Res}\bigl(F(s) e^{st}, s_k\bigr).
    \label{laplace:eq:inversion_residuen}
\end{equation}
Das unendliche Linienintegral ist zu einer endlichen Summe über die
Singularitäten von $F(s)$ geschrumpft.
Die gesamte Information über den zeitlichen Verlauf von $f(t)$ steckt also
in der Lage und Art der Singularitäten der Bildfunktion.

Ein kurzer Test an einem bekannten Korrespondenzpaar zeigt, dass die Formel
hält, was sie verspricht.
Die Bildfunktion $F(s) = \frac{1}{s-a}$ besitzt einen einfachen Pol bei
$s = a$ mit dem Residuum
\begin{equation}
    \operatorname{Res}\left(\frac{e^{st}}{s-a}, a\right)
    =
    \lim_{s \to a} (s-a) \frac{e^{st}}{s-a}
    =
    e^{at}.
    \label{laplace:eq:beispiel_exponential}
\end{equation}
Gleichung~\eqref{laplace:eq:inversion_residuen} liefert also
$f(t) = e^{at}$ und reproduziert damit exakt den Eintrag der
Exponentialfunktion aus Tabelle~\ref{laplace:tab:korrespondenzen}.
Was dort als auswendig gelernte Korrespondenz stand, ist jetzt das Residuum
eines Konturintegrals.

So elegant dieser Mechanismus ist, er trägt eine versteckte Voraussetzung
in sich:
$F(s)$ muss in der ganzen linken Halbebene eine \emph{eindeutige},
meromorphe Funktion sein, denn nur dann dürfen wir die D-Kontur überhaupt
zuziehen.
Genau diese Eindeutigkeit geht verloren, sobald Wurzeln oder Logarithmen
der Variablen $s$ auftreten.
Warum das bei physikalisch völlig harmlosen Problemen passiert und wie die
Kontur dann repariert werden muss, ist das Thema des nächsten Abschnitts.
